% Generated draft: requires mathematical and source review.
\documentclass[11pt]{article}
\usepackage[margin=25mm]{geometry}
\usepackage{fontspec}
\setmainfont{Latin Modern Roman}
\usepackage{amsmath,amssymb,mathtools,unicode-math}
\setmathfont{Latin Modern Math}
\usepackage{enumitem,needspace,xurl,hyperref}
\setlength{\parindent}{0pt}
\setlength{\parskip}{5pt}
\setlength{\emergencystretch}{3em}
\newcommand{\N}{\mathbb{N}}
\newcommand{\Z}{\mathbb{Z}}
\newcommand{\Q}{\mathbb{Q}}
\newcommand{\R}{\mathbb{R}}
\newcommand{\C}{\mathbb{C}}
\newcommand{\F}{\mathbb{F}}
\newcommand{\A}{\mathbb{A}}
\newcommand{\PP}{\mathbb{P}}
\newcommand{\E}{\mathbb{E}}
\newcommand{\eps}{\varepsilon}
\newcommand{\dd}{\,\mathrm{d}}
\begin{document}
\section*{UnsolvedMath ID: 700902}
\textbf{Source problem code:} catalogue entry
\subsection{Definitions and mathematical statement}
For a group \(K\), define its derived series by \(K^{(0)}=K\) and \(K^{(n+1)}=[K^{(n)},K^{(n)}]\) for \(n\geq 0\), where \([L,L]\) is the subgroup generated by all commutators of elements of \(L\). The group \(K\) is soluble if \(K^{(d)}=1\) for some integer \(d\geq 0\); its derived length \(\operatorname{dl}(K)\) is the least such integer. Thus the trivial group has derived length zero.

For a group \(S\), let \(\operatorname{Var}(S)\) denote the variety of groups generated by \(S\), namely the class of all groups satisfying every group identity satisfied by \(S\). Here a group is called two-generated if it has the form \(\langle x,y\rangle\), allowing redundant generators.

The supplied question is whether, for every finite simple group \(S\) subject to the source's qualification ``(known)'', and every integer \(d\geq 0\), the following implication holds:
\[
\bigl(\exists H\leq S\;[H\text{ is soluble and }\operatorname{dl}(H)=d]\bigr)
\ \Longrightarrow\ 
\bigl(\exists G\in\operatorname{Var}(S)\;\exists x,y\in G\;
[G=\langle x,y\rangle,\ G\text{ is soluble},\ \operatorname{dl}(G)=d]\bigr).
\]
The source does not define the class intended by ``(known)'', so the scope of the quantifier over \(S\) remains unresolved. The displayed formulation uses the convention that two-generated means generated by at most two elements; whether the source instead requires minimum generating number exactly two is unspecified. No nonabelian restriction on \(S\), and no finiteness requirement on \(G\), is stated.
\subsection{Short English statement}
Must the variety generated by a finite simple group containing a soluble subgroup of derived length \(d\) also contain a soluble group generated by two elements with derived length exactly \(d\)? The source qualifies the finite simple group as ``(known)'' without defining that restriction.
\subsection{Sources}
\begin{itemize}
\item Supplied local UnsolvedMath registry record: permanent ID \(700902\), exact source problem code \texttt{catalogue entry}, title \emph{Kourovka Notebook Problem 12.85}; upstream revision \texttt{d371ee917ac9647fbb06cd6fc529deec6f3dda82}.
\item Supplied attribution: S. A. Syskin, proposed in 1992. Source: \emph{The Kourovka Notebook (Unsolved Problems in Group Theory)}, 12th issue (1992), Problem 12.85; also identified in the record with the 21st issue edition (2026), arXiv:1401.0300v48.
\item The supplied collection description attributes the notebook to editors E. I. Khukhro and V. D. Mazurov and describes coverage of Issues 1--21 (1965--2026). The registry labels this problem \emph{open}; that label is supplied metadata, not an independently verified literature assessment. No source was independently consulted for this draft.
\end{itemize}
\end{document}
